\lyseis{2.3}{Σύγκριση κλασμάτων}
\sol{1} α) $\frac{5}{9}<\frac{7}{9}$ \quad β) $\frac{4}{5}>\frac{4}{9}$ \quad γ) $\frac{6}{11}<\frac{6}{7}$ \quad δ) $\frac{13}{20}>\frac{9}{20}$ \quad ε) $\frac{1}{3}>\frac{1}{4}$.
\sol{2} α) $\frac{10}{15}>\frac{9}{15}$, άρα $\frac{2}{3}>\frac{3}{5}$ \quad β) $\frac{15}{18}>\frac{14}{18}$, άρα $\frac{5}{6}>\frac{7}{9}$ \quad γ) $\frac{45}{60}>\frac{44}{60}$, άρα $\frac{3}{4}>\frac{11}{15}$ \quad δ) $\frac{14}{24}<\frac{15}{24}$, άρα $\frac{7}{12}<\frac{5}{8}$ \quad ε) $\frac{81}{90}>\frac{80}{90}$, άρα $\frac{9}{10}>\frac{8}{9}$.
\sol{3} Ίδιος αριθμητής: μεγαλύτερο είναι αυτό με τον μικρότερο παρονομαστή. Άρα $\frac{17}{10}>\frac{17}{11}>\frac{17}{12}>\frac{17}{13}>\frac{17}{14}$.
\sol{4} Μικρότερα από το 1: $\frac{7}{9}$, $\frac{4}{21}$, $\frac{99}{100}$. Ίσο με το 1: $\frac{12}{12}$. Μεγαλύτερα από το 1: $\frac{15}{11}$, $\frac{100}{99}$.
\sol{5} Με κοινό παρονομαστή 12: $\frac{9}{12}$, $\frac{10}{12}$, $\frac{8}{12}$, $\frac{7}{12}$, $\frac{11}{12}$. Σε αύξουσα σειρά: $\frac{7}{12}<\frac{2}{3}<\frac{3}{4}<\frac{5}{6}<\frac{11}{12}$.
\sol{6} $\frac{7}{3}$: $2<\frac{7}{3}<3$ (γιατί $2=\frac{6}{3}$). $\frac{11}{4}$: ανάμεσα στο 2 και 3. $\frac{25}{6}$: ανάμεσα στο 4 και 5 ($4\cdot 6=24$). $\frac{40}{9}$: ανάμεσα στο 4 και 5 ($4\cdot 9=36$, $5\cdot 9=45$). $\frac{100}{7}$: ανάμεσα στο 14 και 15 ($14\cdot 7=98$).
\sol{7} Με μονάδα 8 εκ.: το $\frac{1}{4}$ στα 2 εκ., το $\frac{1}{2}$ στα 4 εκ., το $\frac{5}{4}$ στα 10 εκ., το $\frac{3}{2}$ στα 12 εκ. και το $\frac{7}{4}$ στα 14 εκ. από το 0.
\sol{8} Α: $\frac{5}{6}$, Β: $\frac{7}{6}$ (που είναι $1+\frac{1}{6}$), Γ: $\frac{10}{6}$ (ίσο με $\frac{5}{3}$), Δ: $\frac{11}{6}$.
\sol{9} α) Σωστό. β) Λάθος: μεγαλύτερο είναι αυτό με τον μεγαλύτερο αριθμητή. γ) Σωστό. δ) Λάθος: $\frac{5}{5}=1$. ε) Σωστό. στ) Λάθος: $\frac{2}{3}=\frac{10}{15}>\frac{9}{15}=\frac{3}{5}$.
\sol{10} Κοινός παρονομαστής το $\EKP(5,\,9,\,4)=180$: $\frac{4}{5}=\frac{144}{180}$, $\frac{7}{9}=\frac{140}{180}$, $\frac{3}{4}=\frac{135}{180}$. Το πιο επιτυχημένο είναι αυτό με τα $\frac{4}{5}$ και το λιγότερο αυτό με τα $\frac{3}{4}$.
\sol{11} Κοινός παρονομαστής το $\EKP(5,\,7,\,6)=210$: $\frac{3}{5}=\frac{126}{210}$, $\frac{2}{7}=\frac{60}{210}$, $\frac{1}{6}=\frac{35}{210}$. Έτρεξε περισσότερο την πρώτη μέρα, μετά τη δεύτερη και λιγότερο την τρίτη.
\sol{12} $\frac{3}{4}=\frac{15}{20}$ και $\frac{7}{10}=\frac{14}{20}$. Επειδή $15>14$, χρειάζεται περισσότερο από όσο έχει: το πακέτο δεν φτάνει. Λείπει $\frac{1}{20}$ του κιλού.
\sol{13} Κοινός παρονομαστής 24: Πολυξένη $\frac{15}{24}$, Βασίλης $\frac{16}{24}$, Στέλιος $\frac{14}{24}$. Πιο φορτισμένο είναι του Βασίλη, μετά της Πολυξένης και τελευταίο του Στέλιου.
\sol{14} Κοινός παρονομαστής 9: Βασίλης $\frac{4}{9}$, Ζωή $\frac{1}{3}=\frac{3}{9}$, Μυρτώ $\frac{2}{9}$. Σειρά: Βασίλης, Ζωή, Μυρτώ.
\sol{15} $3\cdot 7=21$ και $5\cdot 4=20$, και $21>20$, άρα $\frac{3}{5}>\frac{4}{7}$. Δουλεύει γιατί με κοινό παρονομαστή το $5\cdot 7=35$ έχουμε $\frac{3}{5}=\frac{3\cdot 7}{35}=\frac{21}{35}$ και $\frac{4}{7}=\frac{4\cdot 5}{35}=\frac{20}{35}$: οι αριθμητές που συγκρίνουμε είναι ακριβώς τα χιαστί γινόμενα.
\sol{16} $\frac{2}{5}=\frac{8}{20}$ και $\frac{3}{5}=\frac{12}{20}$. Ανάμεσά τους: $\frac{9}{20}$, $\frac{10}{20}$, $\frac{11}{20}$.
\sol{17} $\frac{1}{3}=\frac{4}{12}$ και $\frac{5}{6}=\frac{10}{12}$. Ανάμεσα: $\frac{5}{12}$, $\frac{6}{12}$, $\frac{7}{12}$, $\frac{8}{12}$, $\frac{9}{12}$, δηλαδή 5 κλάσματα.
\sol{18} Το $\frac{10}{21}$ έχει μεγαλύτερους όρους από το $\frac{2}{3}$, αλλά είναι μικρότερο: $10\cdot 3=30<42=21\cdot 2$. Άρα οι μεγάλοι αριθμοί δεν σημαίνουν μεγάλο κλάσμα.
\sol{19} Όχι. Όσο κοντά κι αν είναι δύο κλάσματα, ανάμεσά τους χωράει άλλο: αν γράψουμε τα δύο κλάσματα με διπλάσιους (ή μεγαλύτερους) όρους, οι αριθμητές απομακρύνονται και χωράει ένας αριθμός ανάμεσα. Άρα δεν υπάρχει «επόμενο» κλάσμα, αντίθετα με τους φυσικούς.
\sol{20} Ενδεικτικά: ύπνος $\frac{8}{24}$, σχολείο $\frac{6}{24}$, διάβασμα $\frac{3}{24}$, παιχνίδι $\frac{3}{24}$, φαγητό και άλλα $\frac{4}{24}$. Με την ίδια ονομασία, συγκρίνεις απλώς τους αριθμητές: ύπνος, σχολείο, φαγητό, και μετά διάβασμα και παιχνίδι ισοδύναμα.
\sol{21} Με μονάδα 10 εκ.: $\frac{1}{5}$ στα 2 εκ., $\frac{3}{10}$ στα 3 εκ., $\frac{1}{2}$ στα 5 εκ., $\frac{7}{10}$ στα 7 εκ., $\frac{4}{5}$ στα 8 εκ., $\frac{6}{5}$ στα 12 εκ., $\frac{3}{2}$ στα 15 εκ. και $\frac{9}{5}$ στα 18 εκ.
